\chapter{奶龙学的对象、边界与元概念}\label{ch:02}

\begin{nlkeypoints}
  \begin{enumerate}[label=(\arabic*),leftmargin=2.4em]
    \item 本书的研究对象是一个\keyw{“正典—变种”结构}：
          官方奶龙（正典 IP）与奶蛙（民间变种）并非两个独立对象，
          而是同一文化过程的两端，奶蛙寄生于奶龙的意义系统而存在。
    \item “奶龙学”不是一个学科，而是一个\keyw{教学标签}：
          它借用学科命名形式，把对象边界、证据分级与分析工具
          打包成可复用的研究单元；本书明确给出它的
          \keyw{不是清单}（不是什么）与\keyw{是清单}（是什么）。
    \item 全书依赖四个元概念：\keyw{meme（迷因）}、
          \keyw{文本盗猎}、\keyw{抽象文化}、\keyw{融合文化}。
          四个概念分别回答“它是什么单位”“它是谁的行为”
          “它属于哪种社群风格”“它发生在什么媒介结构里”。
  \end{enumerate}
\end{nlkeypoints}

绪论交代了“研究什么”与“为什么研究”。本章回答一个更前置的问题：
\keyw{奶龙学研究的对象到底是什么}。这个问题比看上去更困难，
因为对象本身是一个正在移动的靶子：官方在迭代形象，社群在增殖变种，
平台在改写分发规则。对象边界不事先划清楚，后面的每一章
都会把三种不同层面的陈述混在一起。本章先划界，再给出
全书共用的四个元概念。

\section{两个“奶龙”：正典与民间}\label{sec:02:canon}

\begin{definition}[对象结构：正典与变种]\label{def:02:object}
  \textbf{正典奶龙}指由第七印象文化传媒（深圳）有限公司运营的
  原创动画 IP 及其全部官方内容形态（短视频、番剧、有声读物、
  表情包、衍生商品），其形象、角色与叙事由运营方控制并受版权保护
  \cite{jiang2026nailong}。
  \textbf{奶蛙}指从奶龙形象出发、经 AI 重绘发生形体畸变、
  再由社群筛选与再生产而稳定下来的民间变种形象，
  以及与之绑定的笑声、语言与动作系统
  \cite{lostmedia2026naiwa,shoutanji2026naiwa}。
  \textbf{奶龙学}指以“正典—变种”这一对结构及其传播、
  接受与产业化为对象的跨学科教学—研究单元。
\end{definition}

定义里最要紧的不是三个名词，而是那个“对”字。
奶蛙的一切辨识度都来自与正典奶龙的\keyw{相似—差异差值}：
它必须足够像奶龙（否则无法借用奶龙既有流量与情感积累），
又必须足够不像（否则不构成新刺激）。用第六节的语言说，
奶蛙是奶龙的一个\keyw{变异体}（variant），其文化能量
来自“可辨认的偏离”。这也解释了第一版教材的一个观察：
平时人们把奶蛙直接叫作“奶龙”，而奶龙同好会本质上是奶蛙同好会
——称谓的滑移本身就是两个层面相互渗透的证据\conf{C}。

\begin{nltermbox}{术语：正典（canon）}
  借自文学研究：指一个文本传统中被共同体默认为“权威版本”的
  作品集合。在网络二创语境下，“正典”往往不是法律意义上的
  版权主体，而是\keyw{意义权威}：二创者即使不受版权约束，
  仍然会自发地以官方设定为参照系（“奶龙其实没有黑色手”
  这类吐槽本身就是参照系存在的证明）。本书同时使用两个层面：
  法律层面以第七印象为权利主体\cite{jiang2026nailong}，
  意义层面以官方设定为参照系。
\end{nltermbox}

为什么要把“奶蛙”单独立名，而不是继续混用“奶龙”？
三个理由。第一，\keyw{法律}理由：奶蛙二创与正典 IP 之间存在
真实的版权张力（2026 年 3 月官方曾关评、拉黑发奶蛙账号
\conf{C}\cite{shoutanji2026naiwa}；人民日报报道亦显示官方
对“用 AI 工具批量生成低质内容”等侵权形态的维权行动
\cite{jiang2026nailong}），混用名称会模糊权利边界。
第二，\keyw{分析}理由：正典奶龙的成功是\keyw{内容产业}问题
（子供向定位、剧集更新、授权运营），奶蛙的扩散是\keyw{参与文化}
问题（畸变美学、空耳语言、群体仪式），两套解释变量几乎不重叠。
第三，\keyw{教学}理由：学生最容易犯的错误就是把“IP 火”
等同于“梗火”；分名是为了让错误在第一次出现时就被看见。

\section{学科边界：奶龙学不是什么，是什么}\label{sec:02:boundary}

“奶龙学”三个字本身带有戏谑色彩（第一版教材标题
《捧腹大笑思想与奶龙学概论》的“概论”同样是这种风格）。
本书沿用这个命名，但必须立刻声明它的边界，
否则读者会误以为这里存在一个学科。

\begin{table}[htbp]\centering
  \caption{“奶龙学”的边界：不是清单与是清单}\label{tab:02:boundary}
  \small
  \begin{tabularx}{\textwidth}{@{}p{5.2cm}Xp{5.2cm}X@{}}
    \toprule
    & \multicolumn{1}{c}{\heiti 不是什么} & & \multicolumn{1}{c}{\heiti 是什么} \\ \midrule
    对象上 & 不是奶龙动画的艺术史 & & 一个“正典—变种”传播案例 \\
    & 不是第七印象的企业史 & & 迷因文化研究的方法演练场 \\
    & 不是网络热梗盘点 & & 青年情绪与自我呈现的切片 \\ \midrule
    方法上 & 不做经验测量（本书无问卷数据） & & 提供可复用的证据分级制度 \\
    & 不做法律意见（不判断侵权与否） & & 提供可复现的分析流程 \\ \midrule
    伦理上 & 不评价具体个人 & & 以考据伦理为显式教学案例 \\
    & 不背书任何二创行为 & & 以版权张力为分析材料 \\ \bottomrule
  \end{tabularx}
\end{table}

这个边界声明里最容易被忽略的一行是“不做经验测量”。
本书包含大量“分析”，但其中没有任何一项来自本书作者
亲自采集的数据；所有数字要么是 A/B 级公开报道
（如粉丝规模、授权品牌数\cite{jiang2026nailong,chinadaily2026six}），
要么是 C 级社群考据（如起源线索），要么是 M 级教学模型。
本书的价值在于\keyw{把分析做可操作、可复现、可质疑}，
而不在于提供新的经验事实。第十一章给出的问卷
是供修订者与后续研究者施测的\keyw{设计}，不是本书的测量结果。

\section{四个元概念}\label{sec:02:meta}

全书的分析语言由四个元概念支撑。它们分别回答四个问题，
后文各章反复调用，此处只给定义与使用约定。

\subsection{meme（迷因）：最小复制单位}

道金斯把 \en{meme} 定义为文化信息的复制单位：
“就像生物基因在生物体间复制一样，meme 在头脑间复制，
经历变异、选择与保留”\cite{dawkins1976selfish}。
《延伸的表现型》进一步把“复制子”的视角延伸到
文化结构的层面\cite{dawkins1982extended}。
但道金斯式的迷因概念在数字时代面临一个操作难题：
\keyw{复制的边界在哪里}？一条奶蛙视频被转发一万次，
每次转发都算“同一迷因”吗？本书采用 Shifman 的处理方式：
不追问“同一性”的本体论，而是用三维框架
（内容 / 形式 / 立场）给每个用例做\keyw{结构描述}，
再在“结构相似度”高于社群默许阈值时归入同一迷因家族
\cite{shifman2013memes}。第六章将给出奶蛙家族的
归并判据。

\subsection{文本盗猎：行为主体与行为结构}

“文本盗猎者”（textual poachers）指主动挪用、拼贴、
改写既有文本的受众。Jenkins 在《文本盗猎者》一书中
把粉丝写作、混剪与同人活动组织为“参与式文化”的核心形态
\cite{jenkins1992poachers}。奶蛙二创者正是典型盗猎者：
他们盗用正典奶龙的形象资产，嫁接直播连麦的录音资产，
再嫁接 AIGC 工具的生产能力。
本书使用这一概念时的纪律是：\keyw{描述结构，不评价合法性}。
“盗猎”是分析术语，不等于法律上的“侵权”判定，
也不等于价值上的“创意致敬”褒奖；法律层面的张力
在第九章处理。

\subsection{抽象文化：一种社群风格}

“抽象”是中文互联网的原生概念，没有直接对应的外文术语。
本书给出操作性定义：

\begin{definition}[抽象（抽象文化）]\label{def:02:abstract}
  一种以“意义悬置”为核心策略的社群风格：内容不追求
  明确的语义传达，而以“看不懂—被冒犯—仍然好笑”的
  体验回路为组织原则；典型标记包括无逻辑谐音（空耳）、
  故意劣化的制作质量、对严肃语体的戏仿
  （如把宗教文体套用于立牌纪念文本）\conf{C}。
\end{definition}

这个定义是本书构造的教学定义（M 级），它来自对第一版教材
所引材料的归纳（“糖”“倒放大法”“鲱鱼罐头”等社群用语
的使用方式），不是学术共识。定义中“意义悬置”一项最要紧：
它解释了为什么“看不懂”不会杀死传播——
在抽象文化里，\keyw{看不懂本身就是内容}，
是社群内部的准入测试与归属标记
（“看懂了说明你是自己人”）。第十章会看到，
这种“准入测试”功能在高校同好会的纳新中反复出现。

\subsection{融合文化：媒介结构的位置}

Jenkins 用“融合文化”描述内容在多平台间流动、
受众从消费者转为参与者的媒介生态\cite{jenkins2006convergence}；
《可传播的媒体》进一步提出：在融合环境下，
内容的价值不取决于制作方的“黏性”投放，
而取决于受众的“可传播性”行为\cite{jenkins2013spreadable}。
奶蛙恰好卡在这个结构上：正典奶龙由制作方控制分发
（剧集平台、官方账号），奶蛙的“可传播性”则完全由
二创者的转发、二创与线下搬运驱动。
第八章的生产链分析与第九章的产业张力分析，
都以这个“双轨分发”结构为底图。

\section{从梗到“学”：命名的正当性}\label{sec:02:legitimacy}

最后一节处理一个元问题：一本教材叫“××学概论”，
这种命名本身是不是也是一种“梗”？

本书的回答是：\keyw{是，而且这正是它的功能之一}。
“奶龙学概论”的命名模仿了正规学科教材的题录形式，
但它不是伪装——前言已明确本书为开源讲义，
证据分级制度本身就是反伪装的装置。
命名的作用恰恰相反：通过把戏谑对象装进
“最严肃的容器”（教材体例、定义—命题—习题、
GB/T 7714 引文\cite{gbt7714_2015}），它迫使读者
把注意力从“内容好不好笑”转移到“分析是否站得住”。
这是一种\keyw{语域错配}（register mismatch）的自觉使用：
严肃体例 × 戏谑对象，错配本身成为方法立场的表达。

必须同时承认这一策略的限度：教材体例的严肃性
不能替代证据的严肃性。如果本书某处
用 C 级考据支撑了 A 级口吻的结论，
那么错配就从“立场表达”滑向了“误导”——
第十二章的可复现性清单正是为拦截这种滑动而设。

\section{本章小结}\label{sec:02:summary}

本章把研究对象确定为“正典奶龙—民间奶蛙”这对结构，
划定了“奶龙学”的边界（是清单与不是清单各四项），
并给出全书四个元概念：meme（最小复制单位，
操作化为结构相似度）、文本盗猎（行为结构，
不评价合法性）、抽象文化（意义悬置的社群风格，
M 级教学定义）、融合文化（双轨分发结构）。
四个概念分别对应四个问题：单位、行为者、风格、媒介。
第三章起，对象从定义走向历史：
先看正典奶龙如何成为产业（第三章），
再看民间奶蛙的来源考据（第四章）。

\begin{nlexercises}
  \item 给“正典”的\keyw{法律层面}与\keyw{意义层面}
        各举一个在奶蛙语料中的具体实例
        （前者涉及权利主体行为，后者涉及二创者的自发参照），
        并说明两者是否总是一致。
  \item 本章四个元概念中，哪一个最容易被误用为
        “贴标签”（只说“这是文本盗猎”而不展开结构）？
        写出你使用它时必须回答的三个结构问题。
  \item 用定义 2.2 判断下列三个用例是否属于“抽象”：
        (a) 一段把课程作业封面做成“讣告”体的视频；
        (b) 一张画质劣化但语义完整的吐槽表情包；
        (c) 一段完全语义清晰、仅画质劣化的视频。
        逐条写出你的判定依据（意义悬置 / 空耳 / 语体戏仿三项中
        各命中哪几项）。
\end{nlexercises}

\nlkeywords 正典与变种、文本盗猎、抽象文化、融合文化、
meme（迷因）、语域错配
